% ── 6.6 Ablation Study (physical modeling contribution) ──
\subsection{Ablation of Physical Modeling}
\label{sec:ablation}

To quantify the contribution of incidence-angle ($\theta$) and squint-angle
(${\psi_{\text{sq}}}$) modeling to solver performance, we repeat the MOEA-3
experiment with four objective-function variants:
%
\begin{description}
\item[A. Full physical (baseline)] $f_2 = \frac{1}{|\mathcal{S}|}\sum \sin\theta_i \cdot \cos\psi_i$,
      $f_3 = \frac{1}{|\mathcal{S}|}\sum \cos^3\theta_i \cdot \cos^3\psi_i$ (identical to the
      standard configuration in \S~\ref{sec:pareto-coupling}).
\item[B. No squint] $f_2 = \frac{1}{|\mathcal{S}|}\sum \sin\theta_i$,
      $f_3 = \frac{1}{|\mathcal{S}|}\sum \cos^3\theta_i$ (removes ${\psi_{\text{sq}}}$ component).
\item[C. No incidence] $f_2 = \frac{1}{|\mathcal{S}|}\sum \cos\psi_i$,
      $f_3 = \frac{1}{|\mathcal{S}|}\sum \cos^3\psi_i$ (removes $\theta$ component).
\item[D. No physics] $f_2 = f_3 = 1$ (constant per-task objective; solver receives
      no geometric information; $f_2$, $f_3$ shown in Table~\ref{tab:ablation} are post-hoc diagnostics
      computed from the actual observation geometry of the resulting schedule).
      Variant D serves as a null-model control: it tests whether the NSGA-III
      solver's performance is attributable to the physical content of $f_2$
      and $f_3$, or merely to the presence of additional objectives that
      structure the search.
\end{description}
%
All four variants use identical NSGA-III parameters (population $= 100$,
generations $= 200$, Das-Dennis $n=12$), the same 200 scenario files, and
the same G-BL hot-start injection. Only the objective functions differ,
making this a controlled ablation rather than a retuning experiment.
Within-solver ablation with identical NSGA-III parameters and the same
G-BL seed injection controls confounding factors sufficiently that the
Wilcoxon signed-rank test is appropriate for paired comparison of
objective-function variants.
Each variant completes 4 scenario classes $\times$ 50 seeds $= 200$ runs
(total: 800 runs across A--D). Degradation is reported as
$\Delta\% = (\text{A} - \text{V})/\text{A} \times 100$, where A is the full
physical baseline and V is the variant; negative values indicate the variant
outperforms the baseline on $f_1^*$.

\begin{table}[H]
\centering
\caption{Per-class ablation results. Bold: $p < 0.005$ (Wilcoxon signed-rank vs.~A; consistent with the Friedman--Wilcoxon framework of \S\ref{sec:scale-sensitivity}).}
\label{tab:ablation}
\begin{tabular}{llccccrc}
\toprule
\textbf{Class} & \textbf{V} & \textbf{$f_1^*$} & \textbf{$f_2$} &
\textbf{$f_3$} & \textbf{$f_1/n$} & \textbf{$n_\text{sel}$} &
\textbf{$\Delta f_1^*\%$} \\
\midrule
\multirow{4}{*}{S1 ($N=20$)}
 & A & $0.66 \pm 0.28$ & $0.569$ & $0.288$ & $5.93 \pm 1.10$ &
   $11.1 \pm 4.4$ & (baseline) \\
 & B & $0.72 \pm 0.33$ & $0.565$ & $0.281$ & $5.57 \pm 1.05$ &
   $12.8 \pm 5.5$ & $\mathbf{-10.0\%}$ \\
 & C & $0.51 \pm 0.22$ & $0.575$ & $0.300$ & $5.91 \pm 1.24$ &
   $8.8 \pm 3.5$ & $\mathbf{+21.6\%}$ \\
 & D & $0.81 \pm 0.34$ & $0.565$ & $0.266$ & $5.40 \pm 1.04$ &
   $14.7 \pm 5.6$ & $\mathbf{-23.5\%}$ \\
\midrule
\multirow{4}{*}{S2 ($N=100$)}
 & A & $0.86 \pm 0.39$ & $0.561$ & $0.245$ & $5.56 \pm 0.52$ &
   $47.6 \pm 21.4$ & (baseline) \\
 & B & $0.86 \pm 0.39$ & $0.560$ & $0.244$ & $5.55 \pm 0.56$ &
   $47.8 \pm 21.3$ & $+0.0\%$ \\
 & C & $0.83 \pm 0.37$ & $0.561$ & $0.247$ & $5.44 \pm 0.49$ &
   $46.6 \pm 20.0$ & $+3.5\%$ \\
 & D & $0.86 \pm 0.38$ & $0.561$ & $0.243$ & $5.46 \pm 0.60$ &
   $48.6 \pm 21.1$ & $-0.3\%$ \\
\midrule
\multirow{4}{*}{S3 ($N=300$)}
 & A & $0.97 \pm 0.25$ & $0.587$ & $0.175$ & $5.50 \pm 0.23$ &
   $151.1 \pm 38.7$ & (baseline) \\
 & B & $0.98 \pm 0.25$ & $0.587$ & $0.175$ & $5.52 \pm 0.23$ &
   $151.7 \pm 38.5$ & $-0.8\%$ \\
 & C & $0.98 \pm 0.25$ & $0.587$ & $0.175$ & $5.52 \pm 0.23$ &
   $151.1 \pm 38.3$ & $-0.3\%$ \\
 & D & $0.97 \pm 0.25$ & $0.587$ & $0.175$ & $5.51 \pm 0.22$ &
   $151.1 \pm 38.8$ & $-0.2\%$ \\
\midrule
\multirow{4}{*}{S4 ($N=500$)}
 & A & $0.97 \pm 0.28$ & $0.589$ & $0.170$ & $5.51 \pm 0.20$ &
   $190.2 \pm 56.2$ & (baseline) \\
 & B & $0.97 \pm 0.29$ & $0.589$ & $0.169$ & $5.50 \pm 0.19$ &
   $190.6 \pm 56.4$ & $-0.1\%$ \\
 & C & $0.97 \pm 0.28$ & $0.589$ & $0.170$ & $5.51 \pm 0.21$ &
   $191.0 \pm 55.8$ & $-0.4\%$ \\
 & D & $0.96 \pm 0.28$ & $0.589$ & $0.169$ & $5.50 \pm 0.19$ &
   $190.3 \pm 55.6$ & $+0.1\%$ \\
\bottomrule
\end{tabular}
\end{table}

\begin{figure}[H]
    \centering
    \includegraphics[width=\textwidth]{fig6_ablation.pdf}
    \caption{Ablation results across four objective-function variants (A--D) and four scenario classes}
    \label{fig:fig6}
\end{figure}

The ablation results in Table~\ref{tab:ablation} reveal that the influence
of physical objective modeling is concentrated in sparse scenarios and
vanishes as scenario density increases --- a pattern consistent with our
scale-sensitivity analysis (\S~\ref{sec:scale-sensitivity}).

At S1 ($N = 20$), all three ablations produce statistically
significant deviations. Variant C (no incidence) degrades $f_1^*$ by
$+21.6\%$ ($p = 1.9 \times 10^{-5}$), indicating that removing the
incidence-angle component from the objective function forces the solver
into poorer coverage outcomes. Conversely, variant D (no physics)
\emph{improves} apparent $f_1^*$ by $-23.5\%$ ($p = 1.4 \times 10^{-7}$)
by selecting $+32.4\%$ more tasks (14.7 vs.\ 11.1) --- the solver, freed
from geometric quality objectives, converges on raw profit maximization
at the cost of imaging quality. Variant B (no squint) also shows a
significant $-10.0\%$ improvement ($p = 8.0 \times 10^{-5}$),
consistent with the interpretation that removing $\psi_{\text{sq}}$
relaxes the geometric coupling that constrains task selection.

By S2 ($N = 100$), none of the three ablations reaches statistical
significance at $p < 0.005$. The $f_1^*$ of all variants falls within
$\pm 4\%$ of the baseline, and the per-task $f_1$ ($5.44$--$5.56$) and
$n_{\text{sel}}$ ($46.6$--$48.6$) are indistinguishable. The signal from
physical modeling has weakened to the point where it no longer produces
systematic coverage differences.

At S3--S4 ($N \geq 300$), every variant converges to within $1\%$ of the
baseline for $f_1^*$ and to within $0.5\%$ for $f_2$, $f_3$, and $n_{\text{sel}}$.
The per-task quality metrics ($f_1/n$) are identical to two decimal places
($5.50$--$5.52$) across all variants and scene classes. The $f_2$ value
is locked at $0.587$--$0.589$, varying by less than $0.3\%$. This numerical
invariance is the direct signature of the geometric upper bound discussed
in \S~\ref{sec:scale-sensitivity}: once the density of candidate observation
windows saturates the schedule ($N$ exceeds the maximum selectable
targets under the attitude maneuver and squint constraints), the
objective function only penalizes infeasible solutions rather than
differentiating among feasible ones. In this regime, the NSGA-III solver
acts as a constraint-filtered profit maximizer regardless of whether
$f_2$ and $f_3$ carry geometric meaning (A) or serve as constant
placeholders (D).

To verify that the convergence of Variants A and D is not an artifact
of insufficient solver convergence, we ran a control across 5 S3
($N = 300$) and 5 S4 ($N = 500$) scenarios with doubled population
($200$) and generation budget ($400$). Variant A (full physics)
achieves $f_1^* = 0.958 \pm 0.034$ (S3) and $0.892 \pm 0.025$ (S4);
Variant D (no physics) achieves $f_1^* = 0.963 \pm 0.020$ (S3) and
$0.890 \pm 0.026$ (S4), mean differences $\Delta f_1^* = -0.005$ (S3)
and $+0.002$ (S4), with post-hoc $f_2 \approx 0.595$ for both variants. The two variants remain
indistinguishable even with twice the computational budget across all
10 scenarios. Two interpretations are consistent
with this result: (i) the solver convergence to low-squint regimes,
where the empirical correlation $r = 0.93$--$0.98$ exceeds the
definitional null $r_{\text{null}} = -0.51$
(\S~\ref{sec:objective-functions}, Eq.~\ref{eq:rnull}) and aligns
$f_2$ with $f_3$, makes the physical objectives redundant once a
low-squint solution is selected; or (ii) NSGA-III cannot
differentiate between objectives at high density regardless of what
they encode (solver inadequacy). The control favors (i) --- doubling
the budget does not restore differentiation across 10 scenarios
(5 S3 + 5 S4), and a random-init Variant D control (no hot-start,
15 S3 + 15 S4 runs) still converges to low-squint
($f_2 = 0.605 \pm 0.003$, comparable to hot-start D's $0.594$), ruling
out hot-start residue as the sole driver --- but cannot fully rule out
(ii) without exact methods (documented intractable at $N \geq 20$ in
\S~\ref{sec:runtime}). We
therefore state the finding boundedly: under the tested
configuration, removing physical objectives produces
indistinguishable results, and the more parsimonious explanation is
solver convergence to a low-squint regime where the objectives
co-vary rather than mathematical irrelevance of the physical
objects.
